\documentclass{beamer}
% \includeonlyframes{uj}
\synctex=1
\usepackage[utf8]{inputenc}
\usepackage[ngerman]{babel}
\usepackage[T1]{fontenc}
\usepackage{graphicx}
\usepackage[absolute,overlay]{textpos}
\usepackage{csquotes}
\usepackage{eurosym}
\usepackage{wasysym}
\MakeAutoQuote{„}{”}
\synctex=1
\usepackage{lmodern}


% \usepackage{lmodern}
% \usepackage{anyfontsize}

\hypersetup{
  colorlinks=true,
  linkcolor=[rgb]{1, 1, 1}, %
  urlcolor=[rgb]{.2, .2, .5} %
  % allcolors=[rgb]{0, 0, 0} % schwarz
}

\usetheme{Dresden}
\setbeamertemplate{navigation symbols}{}
\setbeamertemplate{footline}[frame number]

\title{Verschlüsselte Kommunikation: Warum und wie.}
\subtitle{Kreta WiSe 2017/18}

% 
% Logo rechts oben:
% \addtobeamertemplate{frametitle}{}{%
%   \begin{textblock*}{130mm}(.963\textwidth,8.4mm)
%     \includegraphics[width=2.25cm]{img-src/fsfw-logo.pdf}
%   \end{textblock*}
% }

\begin{document}

\begin{frame}
  \begin{center}%
  [Kreta-Logo] + 
    \includegraphics[width=3cm]{img-src/fsfw-logo-with-text.pdf}\\[4mm]
    \vspace*{-0.5\baselineskip}

    \parbox{.95\columnwidth}{\centering\Large\inserttitle}

    \vspace*{\baselineskip}

    \structure{\large \insertsubtitle}
  \end{center}
\end{frame}


%%%%%%%%%%%%%%%%%%%%%%%%%%%%%%%%%%%%%%%%%%%%%%%%%%%%%%%%%%%%%%%%%%%%%%%%%%%%%%%%

\begin{frame}[label=ct1]
  \frametitle{Wer sind wir? $\rightarrow$ FSFW}

  \begin{columns}
 
\column[t]{1\textwidth}
\vspace{-5mm}
%   \onslide<+->

  \begin{itemize}
  \item Hochschulgruppe seit 2014, ca. 10 Leute (TU, HTW, ...)
%   \item Hochschulen als Zielgruppe (Multiplikationswirkung)\\
%     und Arbeitsfeld (Räume, Strukturen)
  \item Warum machen wir das? Aus Überzeugung!

%     \bigskip\onslide<+->

  \begin{itemize}
  \item \emph{Überzeugung 1}: freie und quelloffene Software ist (oft) besser\\
    (technische + nicht technische Argumente)\\
    \bigskip
    \pause
  \item \emph{Überzeugung 2}: öffentlich finanzierte wissenschaftliche Inhalte
    (AutorInnen, GutachterInnen) sollten nicht von öffentlich finanzierten
    Bibliotheken für horrende Summen von Zeitschriften-Verlagen gekauft werden
    müssen
  \end{itemize}

    \pause
  \item Bisherige Projekte
    \begin{itemize}
    \item Linux-Install-Party, Linux-Presentation-Day
    \item Monatliche Sprechstunde zu \LaTeX{} u.a.
    \item Programmpapier
    \item \textbf<4-5>{„Uni-Stick”:~80 $\times$ 8\,GB mit freier Software}
    \item Verschlüsselungsgewinnspiel
    \end{itemize}
    \pause
    \pause
    \pause
    \item Für (Mitmachen-)Interessierte: \url{https://fsfw-dresden.de/}
  \end{itemize}
  
\begin{textblock*}{5cm}[0.,0.](90mm,63mm)
\only<4>{
\includegraphics[width=35mm]{img-src/usb-hub}
}
\only<5->{
\includegraphics[width=35mm]{img-src/uni-stick-ausgabe-vortrag}
}
\end{textblock*}  

\column[t]{0.1\textwidth}
~
\end{columns}
  
\end{frame}

%%%%%%%%%%%%%%%%%%%%%%%%%%%%%%%%%%%%%%%%%%%%%%%%%%%%%%%%%%%%%%%%%%%%%%%%%%%%%%%%

\begin{frame}[label=ol1]
  \frametitle{Gliederung}
  \begin{itemize}
   \item Einführung
   \item Sammlung von Argumenten und Diskussion\\
   ("`Ich habe doch nichts zu verbergen"')
   \item Theoretischer Hintergrund und Begriffsklärungen
   \item Unterstützung bei Installation und Einrichtung (Zeit?)
   \begin{itemize}
    \item GPG-Enigmail (Thunderbird)
   \item Mailvelope (Webmail)
   \end{itemize}

  \end{itemize}

\end{frame}

%%%%%%%%%%%%%%%%%%%%%%%%%%%%%%%%%%%%%%%%%%%%%%%%%%%%%%%%%%%%%%%%%%%%%%%%%%%%%%%%

\begin{frame}[label=ef1]
  \frametitle{Einführung }
  In der Geschichte der Menschheit...
  \begin{itemize}
   \item  ... kein Mangel an Beispielen für \textbf{rücksichtslosen Egoismus}
   \item Ausnutzung von Macht zu persönlichem Vorteil
   \pause
   \item "`Wissen ist Macht"'
   
   \item[$\Rightarrow$] sensibler Umgang mit Informationen empfehlenswert\\[2mm]
   
   \pause
   \item Digitalisierung verstärkt das Problem:
   \item E-Mail${}^1$ ist wie Postkarte: unterwegs${}^{2}$ lesbar
   \item E-Mail ist schlimmer als Postkarte:
   \begin{itemize}
    \item automatisiert auswertbar
    \item unbemerkt kopierbar
    \item unbemerkt veränderbar (inkl. Absender)
   \end{itemize}

  \end{itemize}
  
  ~\\[5mm]
  {\tiny $1$: E-Mail = Beispielmedium \qquad $2$: ggf. um die ganze Welt }

\end{frame}

%%%%%%%%%%%%%%%%%%%%%%%%%%%%%%%%%%%%%%%%%%%%%%%%%%%%%%%%%%%%%%%%%%%%%%%%%%%%%%%%

\begin{frame}[label=ef2]

  \frametitle{Einführung 2}

Schutzziele sicherer Kommunikation:
  \begin{itemize}
   \item[$\square$] Vertraulichkeit
   
   \item[$\square$] Integrität (keine Veränderung)
   
   \item[$\square$] Anonymität
  \end{itemize}

\end{frame}
 
%%%%%%%%%%%%%%%%%%%%%%%%%%%%%%%%%%%%%%%%%%%%%%%%%%%%%%%%%%%%%%%%%%%%%%%%%%%%%%%

\begin{frame}[label=d1]
  \frametitle{"`Ich habe doch nichts zu verbergen"'}
  10min freie Diskussion
  
  \begin{itemize}
   \item  je 2-5 Personen
   \item verschiedene Perspektiven\\
  (Bürgerrechtler*in, Innenminister*in)
   \item Argumente aufschreiben
  \end{itemize}
  \pause
  Stichworte:
  \begin{itemize}
   \item Kriminalität (Einbruch, Erpressung, ...)
   \item Privatsphäre
   \item Selbstzensur (analog: Kamera-Attrappen)
   \item Demokratie
   \item Journalismus
   \item Whistleblowing
   \item Erstarken totalitärer Strukturen
   
  \end{itemize}

  
\end{frame}




%%%%%%%%%%%%%%%%%%%%%%%%%%%%%%%%%%%%%%%%%%%%%%%%%%%%%%%%%%%%%%%%%%%%%%%%%%%%%%%

\begin{frame}[label=begriffe]
  \frametitle{Begriffe}
  
  \begin{columns}
\column[t]{0.48\textwidth}

  \begin{itemize}
   \item asymmetrisches Verfahren
   \item Privater Schlüssel
   \item Öffentlicher Schlüssel
   \item Verschlüsseln
   \item Signieren / Signatur
   \item Schlüsselserver
   \item Web of Trust
   \item Widerrufszertifikat
   \item Cryptoparty
  \end{itemize}
  \column[t]{0.6\textwidth}
  \begin{itemize}
   \item Fingerabdruck
   \item Metadaten
   \item GPG vs PGP {\tiny vs S/MIME}
   \item Enigmail vs Mailvelope
  \end{itemize}
\end{columns}
\end{frame}
%%%%%%%%%%%%%%%%%%%%%%%%%%%%%%%%%%%%%%%%%%%%%%%%%%%%%%%%%%%%%%%%%%%%%%%%%%%%%%%%


\begin{frame}[label=bg2]
  \frametitle{Hintergrundwissen 1}
  
  \begin{itemize}
   \item Schlüsselverteilungsproblem
   \item[$\Rightarrow$] asymmetrisches Verschlüsselungsverfahren
   \begin{itemize}
    \item $\exists$ \textbf{ö}ffentlicher \textbf{S}chlüssel
    und \textbf{p}rivater \textbf{S}chlüssel
    \item sehr große zufällige Zahlen (dargestellt als \href{https://cknoll.github.io/files/pub-key-carsten.txt}{Zeichen}%
    \href{file:///home/ck/pub-key-carsten-demo.txt}{salat})
   \end{itemize}
   \pause
   % padlock: <div>Icons made by <a href="http://www.flaticon.com/authors/dave-gandy" title="Dave Gandy">Dave Gandy</a> from <a href="http://www.flaticon.com" title="Flaticon">www.flaticon.com</a> is licensed by <a href="http://creativecommons.org/licenses/by/3.0/" title="Creative Commons BY 3.0" target="_blank">CC 3.0 BY</a></div>
   \item ÖS: zum Verschlüsseln \hspace{5mm}\includegraphics[width=10mm]{img-src/padlock-lock}
   
   %key: <div>Icons made by <a href="http://www.flaticon.com/authors/yannick" title="Yannick">Yannick</a> from <a href="http://www.flaticon.com" title="Flaticon">www.flaticon.com</a> is licensed by <a href="http://creativecommons.org/licenses/by/3.0/" title="Creative Commons BY 3.0" target="_blank">CC 3.0 BY</a></div>

   
   \item PS: zum Entschlüsseln \hspace{5mm}{\vspace*{4mm}\includegraphics[width=19mm]{img-src/padlock-unlock-with-key}}

  \end{itemize}
\pause
  Oft eingesetzt: \texttt{GPG} (\textbf{G}nu\textbf{P}rivacy\textbf{G}uard) $=$ Freie Software\\[2mm]
  {\tiny Vorläufer: \texttt{PGP} (\textbf{P}retty \textbf{G}ood \textbf{P}rivacy) $\neq$ Freie Software, Namensgeber des Verfahrens}\\
  {\tiny S/MIME: anderes Verfahren, hier nicht weiter behandelt}
  
  
\end{frame}

%%%%%%%%%%%%%%%%%%%%%%%%%%%%%%%%%%%%%%%%%%%%%%%%%%%%%%%%%%%%%%%%%%%%%%%%%%%%%%%

\begin{frame}[label=bg3]
  \frametitle{Hintergrundwissen 2}
  
    \begin{itemize}
   % padlock: <div>Icons made by <a href="http://www.flaticon.com/authors/dave-gandy" title="Dave Gandy">Dave Gandy</a> from <a href="http://www.flaticon.com" title="Flaticon">www.flaticon.com</a> is licensed by <a href="http://creativecommons.org/licenses/by/3.0/" title="Creative Commons BY 3.0" target="_blank">CC 3.0 BY</a></div>
   \item Öffentlicher Schlüssel ("`public key"')
   \begin{itemize}
    \item Benötigt zum Verschlüsseln
    \item Sollten alle haben, von denen man verschlüsselte Mails empfangen möchte
    \item kann/sollte man weitergeben $\rightarrow$ auf Keyserver hochladen\\
    {\tiny Bsp: \url{http://pgp.mit.edu} \quad Bedenkenswert: Nicht löschen, nur Widerrufen $\rightarrow$ Anonymität?}
    
   \end{itemize}
   \pause
   \item Privater Schlüssel ("`private key"')
   \begin{itemize}
    \item Benötigt zum Entschlüsseln
    \item Darf nicht verloren gehen\\$\rightarrow$ Entschlüsseln wäre dann unmöglich\\[1mm]
    \item Darf nicht in fremde Hände kommen\\ $\rightarrow$ andere können meine Mails entschlüsseln\\[1mm]
    \item Typischerweise nochmal zusätzlich mit einem Passwort verschlüsselt\\[2mm]
    \pause
    \item[$\Rightarrow$] Sicheres Backup wichtig
   \end{itemize}


  \end{itemize}
  
\end{frame}

\begin{frame}[label=bg3a]{Hintergrundwissen 3}
Was wird verschlüsselt?
\begin{itemize}
\item Text
\item Anhänge {\tiny Format-Empfehlung: \texttt{PGP/MIME} \quad nicht: \texttt{Inline PGP}}
\item[]
\end{itemize}
Was wird nicht Verschlüsselt?
\begin{itemize}
 \item Metadaten
 \begin{itemize}
  \item Absender*in, Empfänger*in, \textbf{Betreff!}, ...
  \item[]
 \end{itemize}
\end{itemize}
Umsetzung:
\begin{itemize}
 \item Enigmail (GPG-Plugin für Thunderbird)
 \item Mailvelope (GPG-Plugin für Firefox / Chrome $\rightarrow$ Webmail)
\end{itemize}


\end{frame}

%%%%%%%%%%%%%%%%%%%%%%%%%%%%%%%%%%%%%%%%%%%%%%%%%%%%%%%%%%%%%%%%%%%%%%%%%%%%%%%%

%%%%%%%%%%%%%%%%%%%%%%%%%%%%%%%%%%%%%%%%%%%%%%%%%%%%%%%%%%%%%%%%%%%%%%%%%%%%%%%

\begin{frame}[label=bg4]{Hintergrundwissen 4}


Schutzziele sicherer Kommunikation:
  \begin{itemize}
   \item[$\square$\visible<2->{\hspace{-.67em}\raisebox{0.1em}{\scalebox{1.3}{\checkmark}}}]
   Vertraulichkeit $\rightarrow$ Verschlüsselung 
   
   \item[$\square$\visible<5->{\hspace{-.67em}\raisebox{0.1em}{\scalebox{1.3}{\checkmark}}}]
   Integrität (keine Veränderung) 
   
   \item[$\square$\scalebox{1.25}{~}] Anonymität
  \end{itemize}
  \pause
  \pause
  ~\\[3mm]
  
  
  Integrität durch Signieren
  
  \begin{itemize}
   \item Prinzip: mit PS verschlüsselte Prüfsumme der Nachricht
   \item Überprüfen $=$ Entschlüsseln mit ÖS
   \item[$\Rightarrow$] Man muss dem öffentlichen Schlüssel vertrauen\\[2mm]
   \pause
   \item Schlüssel signieren bei persönlichem Treffen $\rightarrow$ "`Web of trust"'
   \item Fingerabdruck-Bsp: {\tiny \texttt{214E 4E9D B193 6AF2 CFDC 68DF 13AD 3604 9D3E F6BF}}
  \end{itemize}

  
\end{frame}

%%%%%%%%%%%%%%%%%%%%%%%%%%%%%%%%%%%%%%%%%%%%%%%%%%%%%%%%%%%%%%%%%%%%%%%%%%%%%%%%

\begin{frame}[label=bg5]{Hintergrundwissen 5}
$\square$\visible<3->{\hspace{-.67em}\raisebox{0.1em}{\scalebox{1.3}{\checkmark}}} Anonymität mittels Tor (Ehem. \textbf{T}he \textbf{O}nion \textbf{R}outer)

\begin{center}
\includegraphics[width=0.6\columnwidth]{img-src/TOR_Arbeitsweise}
\end{center}

\pause 
\begin{itemize}
\item {\small \url{https://www.torproject.org/projects/torbrowser.html}}
\item Nachteile: Geschwindigkeit
\end{itemize}
\end{frame}

%%%%%%%%%%%%%%%%%%%%%%%%%%%%%%%%%%%%%%%%%%%%%%%%%%%%%%%%%%%%%%%%%%%%%%%%%%%%%%%%


\againframe{begriffe}



%%%%%%%%%%%%%%%%%%%%%%%%%%%%%%%%%%%%%%%%%%%%%%%%%%%%%%%%%%%%%%%%%%%%%%%%%%%%%%%%
\begin{frame}[label=lk1]{Links}
\begin{itemize}
\item[] \hspace{-2em} PGP-Verschlüsselung in Thunderbird (Plugin "`Enigmail"')
\item \url{https://fsfw-dresden.de/gpg}\\[4mm]
\item[] \hspace{-2em} PGP-Verschlüsselung für Webmail
\item \url{https://vimeo.com/178702500} (Screencast)\\[2mm]
\item \href{https://posteo.de/hilfe/wie-installiere-ich-eine-ende-zu-ende-verschluesselung-pgp-im-browser}{Anleitung von Posteo}
\item[] \hspace{-2em} Allgemeine Infos
\item \url{https://virtual-privacy.org/}
\item[] \hspace{-2em} Cryptopartys
\item {\tiny \url{https://www.cryptoparty.in/}; \quad | \quad \url{https://de.wikipedia.org/wiki/CryptoParty}}
\item[] \hspace{-2em} Videos
\item 
{\tiny \href{https://www.ted.com/talks/glenn_greenwald_why_privacy_matters}{G. Greenwald: \textit{Why privacy matters}}\quad | \quad
\href{https://www.youtube.com/watch?v=XEVlyP4_11M}{J. Oliver + E. Snwoden (lustig)}}
\end{itemize}
\end{frame}

%%%%%%%%%%%%%%%%%%%%%%%%%%%%%%%%%%%%%%%%%%%%%%%%%%%%%%%%%%%%%%%%%%%%%%%%%%%%%%%%
\begin{frame}[label=uj]{Und jetzt?}
\begin{itemize}
 \item Fragen?\\[10mm]
 \item Unterstützung jetzt bei Installation
 \begin{itemize}
  \item Enigmail (Plugin für Thunderbird)
  \item Mailvelope (Browserplugin für Webmail)
  \item \scalebox{1.9}{$\triangle$}\hspace{-1.0em}\raisebox{.2em}{!} \hspace{0.5em} Zeit?\\[5mm]
  \pause
  \end{itemize}
  \item Unterstützung später {\tiny (im Rahmen unserer Möglichkeiten)}:
  \begin{itemize}
   \item \url{https://fsfw-dresden.de/gpg} (diese Vortragsfolien)
   \item \url{https://fsfw-dresden.de/sprechstunde}
   \item \url{kontakt@fsfw-dresden.de}
  
 \end{itemize}
 \pause
 \item[]
 \item Verschlüsselung regelmäßig nutzen, Erfahrung sammeln
 \item Selber aktiv werden {\tiny "`Wer, wenn nicht wir? Wann, wenn nicht jetzt?"'}

\end{itemize}



\end{frame}

%%%%%%%%%%%%%%%%%%%%%%%%%%%%%%%%%%%%%%%%%%%%%%%%%%%%%%%%%%%%%%%%%%%%%%%%%%%%%%%%
\end{document}
